\section{\blue{Conclusion}}
\label{sec:conclusion}

\blue{
Forecasting models have grown ever more elaborate, and this article asked whether periodic sensor streams require it. They do not. A window of a periodic series carries three kinds of information, the shape, the level and the amplitude of its cycles, which form one low-dimensional structure, a level term plus an amplitude-scaled shape that captures 78\% to 93\% of a normalized window on seven benchmarks. \sysplus gives each part a token and a small pathway: phase tokens, which stay in a fixed subspace when every cycle is reshaped in the same way, carry the shape through a routing Transformer with a few routers and one shared predictor; a level token, the cycle means, is mapped onto three fixed horizon profiles; and an amplitude token, computed from the level, scales the cycle shape while provably leaving every future cycle mean unchanged. A unified low-dimensional view shows why the three tokens are needed and why each pathway takes its form: for fixed weights, a forecaster that reads phase tokens through a layer normalization confines its forecast level and amplitude to two envelopes that no added depth removes, and what lies beyond them is a few profile coefficients and one level-driven factor. Read from the trained models, the envelopes bind exactly where the level and amplitude tokens gain most.

\sysplus attains the best average MSE against eight baselines on six of the seven benchmarks with about 1.2K parameters in its phase pathway and at most about 2.1K more for the level and the amplitude on the largest benchmark, Traffic, where it cuts the FLOPs of PatchTST and Crossformer by about 99.99\%. Because one shared model serves every sensor, these savings recur for every sensor of a network at every update, and the level--amplitude coupling that the amplitude token carries is strongest on the two sensor networks, Traffic and Electricity. More broadly, the results suggest that the capacity a forecaster needs is set by what its tokens leave entangled: with tokens that cover what the forecast needs, a small model suffices.

\fakeparagraph{Limitations and future work}
\sysplus is channel-independent: it exploits the scale of a sensor network through parallelism and one shared model, and combining its tokens with explicit spatial modeling is a natural next step. The capability parts of our results are existence results under the seasonal random walk and its multiplicative variant, whereas the envelopes hold for every parameter value and are measured on the trained models. The analysis assumes a single dominant period and horizons that are whole numbers of cycles; multiple interacting periods and series without a clear period are left to future work, as are input-dependent gates between the pathways.
}
